\chapter{Arquitectura C4}
\label{cap:arquitectura-c4}

Este cap\'itulo integra la documentaci\'on arquitect\'onica ya versionada en
\texttt{docs/diagrams/} (Nivel 2 y Nivel 3 del modelo C4~\cite{c4-model}) y las decisiones
de arquitectura relevantes (\texttt{docs/adr/}). Las fuentes de diagrama
son Graphviz (\texttt{.dot}) y PlantUML (\texttt{.puml}) planos, sin macros
de C4-PlantUML, siguiendo la convenci\'on ya establecida en el repositorio.

\section{C4 Nivel 2 --- Contenedores}
\label{sec:c4-nivel-2}

El diagrama de contenedores (\texttt{docs/diagrams/c4-contenedores/})
describe los cuatro contenedores del sistema: el \emph{frontend} Angular, el
\emph{backend} Spring Boot, PostgreSQL y Redis, junto con los protocolos
reales entre ellos (HTTPS/JSON con cookies entre frontend y backend; JDBC
entre backend y PostgreSQL; el cliente Redis de Spring Data entre backend y
Redis). Su fuente ya cuenta con una imagen renderizada en el repositorio
(\texttt{docs/diagrams/c4-contenedores/c4-contenedores.png}).

% EVIDENCIA-COMPARTIDA-C4-L2
% Ruta esperada: figuras/compartidas/c4-contenedores.png
\evidencia{figuras/compartidas/c4-contenedores.png}%
{Diagrama C4 Nivel 2 (contenedores) de BIOPET: frontend, backend, PostgreSQL y Redis, con los protocolos reales entre ellos.}%
{fig:c4-nivel-2}

\section{C4 Nivel 3 --- Componentes del backend}
\label{sec:c4-nivel-3}

El diagrama de componentes del backend
(\texttt{docs/diagrams/c4-componentes-backend/}) fue elaborado mediante
inspecci\'on directa del c\'odigo fuente (constructores, llamadas
directas, configuraci\'on de Spring), no por convenci\'on de nombres.
Documenta 20~componentes reales agrupados en seis \'areas y 4~elementos
externos:

\begin{itemize}[nosep]
  \item \textbf{API REST / Controladores:} \texttt{AuthController},
    \texttt{MascotaController}, \texttt{UsuarioController}.
  \item \textbf{Servicios de aplicaci\'on:} \texttt{AuthService},
    \texttt{MascotaService}.
  \item \textbf{Seguridad y autenticaci\'on:} \texttt{SecurityConfig},
    \texttt{JwtAuthenticationFilter}, \texttt{JwtService},
    \texttt{JwtCookieService}, \texttt{UserDetailsServiceImpl},
    \texttt{LoginRateLimiterService}, \texttt{TokenBlacklistService},
    \texttt{AuthenticationAuditService},
    \texttt{ProblemAuthenticationEntryPoint},
    \texttt{ProblemAccessDeniedHandler}.
  \item \textbf{Manejo de errores:} \texttt{GlobalExceptionHandler},
    \texttt{ProblemDetailFactory}.
  \item \textbf{Persistencia:} \texttt{UsuarioRepository},
    \texttt{MascotaRepository}.
  \item \textbf{Infraestructura del servidor:} \texttt{TomcatDualConnectorConfig}.
  \item \textbf{Elementos externos:} Frontend BIOPET, PostgreSQL, Redis,
    logs locales del proceso.
\end{itemize}

Este diagrama no es un diagrama de clases: no representa atributos,
m\'etodos ni constructores; los DTO y entidades se mencionan \'unicamente
en prosa. Su documento completo, con las siete relaciones principales
(login, autenticaci\'on de una solicitud protegida, consulta/modificaci\'on
de mascotas, logout y revocaci\'on, rate limiting, manejo de errores v\'ia
\texttt{ProblemDetail} y registro de auditor\'ia), est\'a en
\texttt{docs/diagrams/c4-componentes-backend/C4-L3-backend.md}.

\textbf{Nota sobre el renderizado:} la fuente \texttt{c4-componentes-backend.dot}
ya cuenta con una imagen \texttt{.png} generada con Graphviz 15.1.0
(\texttt{dot -Tpng c4-componentes-backend.dot -o c4-componentes-backend.png},
sintaxis validada previamente con \texttt{dot -Tcanon}), disponible en
\texttt{docs/diagrams/c4-componentes-backend/c4-componentes-backend.png}
(5594\texttimes{}2190~px). Falta copiar esa imagen a
\texttt{figuras/jaime/c4-componentes-backend.png} para que aparezca
autom\'aticamente en este informe.

% EVIDENCIA-JAIME-C4-L3
% Ruta esperada: figuras/jaime/c4-componentes-backend.png
\evidencia{figuras/jaime/c4-componentes-backend.png}%
{Diagrama C4 Nivel 3 (componentes) del backend de BIOPET: controladores, servicios, seguridad, manejo de errores, persistencia e infraestructura del servidor.}%
{fig:c4-nivel-3}

\section{Decisiones de arquitectura relevantes (ADR)}
\label{sec:adr-relevantes}

\begin{table}[H]
  \centering
  \caption{ADR vigentes del repositorio y su relaci\'on con la arquitectura descrita.}
  \label{tab:adr-relevantes}
  \begin{tabular}{@{}p{2cm}p{9.5cm}p{2.5cm}@{}}
    \toprule
    ADR & Decisi\'on & Estado \\
    \midrule
    ADR-002 & Migraci\'on de la pila de backend de ASP.NET Core~8 a Java~21 / Spring Boot~3.2. & Aceptado \\
    ADR-003 & Revocaci\'on de JWT mediante lista negra en Redis, indexada por \texttt{jti}. & Aceptado \\
    ADR-004 & Estrategia de base de datos reproducible: Flyway como fuente de verdad del esquema, m\'as \texttt{db/schema.sql}/\texttt{db/roles.sql}/\texttt{db/seed.sql} para inicializaci\'on de Docker con privilegios m\'inimos. & Aceptado \\
    ADR-005 & Reproducibilidad del despliegue: \texttt{Makefile} y pinning de im\'agenes por digest \texttt{sha256}. & Aceptado \\
    ADR-006 & Decisiones de autenticaci\'on y seguridad: JWT por cookies, claims, revocaci\'on, rate limiting, cabeceras, TLS. & Aceptado \\
    \bottomrule
  \end{tabular}
\end{table}

ADR-001 (selecci\'on inicial de ASP.NET Core~8 en la Entrega~1A) qued\'o
formalmente superado por ADR-002 y se conserva \'unicamente como registro
hist\'orico fuera de este repositorio de la Tercera Entrega.
